\chapter{Build e Deploy}

\section{Pipeline de empacotamento}

O Void Descent é distribuído como executável \texttt{.exe} portátil para Windows
x64, empacotado via Electron e \texttt{electron-builder}. Para representar a
topologia de \emph{nodos} físicos e lógicos, e os \emph{artifacts} que circulam
entre eles, o tipo de diagrama UML adequado é o \textbf{Deployment Diagram}.

\diagfig{diagrams/deployment.png}{Diagrama de Deployment do \emph{pipeline} de \emph{build} e \emph{runtime}}{deployment}

\section{Comandos de \emph{build}}

\subsection{Desenvolvimento local}

Para rodar a versão de desenvolvimento, com \emph{hot reload} manual via
\emph{refresh} do Electron:

\begin{verbatim}
npm install
npm start
\end{verbatim}

\subsection{Build do \texttt{.exe} para Windows}

A configuração em \texttt{package.json} define o \emph{target} \texttt{portable},
um executável auto-extraível sem instalador:

\begin{verbatim}
npm run build:win
\end{verbatim}

A saída é \texttt{dist/VoidDescent-1.0.0.exe}, com $\sim$73\,MB. Em paralelo,
\texttt{electron-builder} produz \texttt{dist/win-unpacked/} (a árvore de arquivos
antes do empacotamento, útil para \emph{debug}) e arquivos auxiliares
(\texttt{builder-debug.yml}, \emph{blockmaps}) que são ignorados pelo
\texttt{.gitignore}.

\section{Cross-compile do macOS}

A geração do \texttt{.exe} para Windows acontece no \emph{host} macOS, sem
necessidade de Wine ou máquina virtual Windows. O \emph{target} \texttt{portable} é
uma das poucas configurações que o \texttt{electron-builder} cross-compila
nativamente; alvos NSIS (instalador) ou MSI exigem Wine no \texttt{PATH}.

\section{Por que 73\,MB?}

A composição aproximada do binário \texttt{.exe} é a seguinte:

\begin{tabular}{lr}
  \toprule
  \textbf{Componente} & \textbf{Tamanho} \\
  \midrule
  Chromium (V8 e Blink) & $\sim$60\,MB \\
  Node.js \emph{runtime} & $\sim$10\,MB \\
  Locales (ICU) & $\sim$5\,MB \\
  Glue do Electron & $\sim$3\,MB \\
  Bundle do jogo (\texttt{app.asar}) & $\sim$80\,KB \\
  \bottomrule
\end{tabular}

O ``jogo em si'' ocupa $0{,}11\%$ do binário. Os outros $99{,}89\%$ são o
\emph{runtime} necessário para hospedar uma página WebGL como aplicação
\emph{desktop}. Esta foi a escolha explícita do GDD; a alternativa seria usar o
\emph{webview} do sistema (Tauri ou Neutralino reduziriam o tamanho para 2 a
10\,MB), mas o requisito de portabilidade absoluta e reprodutibilidade pesou na
decisão por Electron.

\section{Distribuição}

O \texttt{.exe} é \emph{não-assinado}, pois não há certificado de \emph{code
signing} de desenvolvedor Microsoft. Na primeira execução em Windows 10 ou 11, o
\emph{SmartScreen} exibe o aviso ``O Windows protegeu seu PC''. A interação
esperada do usuário é:

\begin{enumerate}
  \item Clicar em \textit{Mais informações}.
  \item Clicar em \textit{Executar mesmo assim}.
\end{enumerate}

Após essa primeira execução, o \emph{hash} do executável é registrado localmente e
o aviso não reaparece. O modelo é coerente com a distribuição via \emph{releases}
do GitHub, o canal escolhido para entrega.

\section{Repositório}

O código-fonte e o \texttt{.exe} compilado estão disponíveis em:

\begin{center}
  \url{https://github.com/luisg0c/Void-Descent}
\end{center}

A história de \emph{commits} reflete a progressão real do projeto entre 03 de
março e 09 de maio de 2026, com mensagens em formato \emph{Conventional Commits}.
